% Part: normal-modal-logic
% Chapter: axioms-systems
% Section: introduction

\documentclass[../../../include/open-logic-section]{subfiles}

\begin{document}

\olfileid{nml}{prf}{int}

\olsection{مقدمة}

قد فرغنا من دلالة اللغة الجهية الأساسية بالنماذج، ومن معنى صلاحية
!!a{formula}: إما أن تصدق في كل عالم من كل نموذج، وإما أن يختص هذا
التعميم بفئة من النماذج أو بالأطر والنماذج المبنية عليها.
ونريد الآن وصل هذا الوصف الدلالي بمعنى برهاني، هو
!!{derivability}. فالمطلوب حدّ !!{derivability} في نسق يجعل
!!a{formula} !!{derivable} متى وفقط متى كانت صالحة.

وأقدم نظم !!{derivation} وأبسطها النظم المسماة هلبرتية أو بديهية.
ويسهل نسبيًا إنشاء نظم !!{derivation} من هذا النمط لكثير من
المنطقات الجهية؛ فنظم !!{derivation} الهلبرتية يسيرة من جهة
دراستها الميتانظرية، كبحث السلامة والتمام. ولا يلزم من يسر الدراسة
يسر الاستعمال: فإثبات !!{formula}s فيها أشق بكثير من إثباتها
في الاستنتاج الطبيعي، مثلًا.

في نظم !!{derivation} الهلبرتية، !!a{derivation} لـ!!a{formula}
متتالية من !!{formula}s تبتدئ ببديهيات وتنتهي، بقواعد استدلال
قليلة، إلى ال!!{formula} المطلوبة. ولكي يوافق نسق !!{derivation}
الدلالة، لا بد أن تكون الصيغ !!{derivable} صادقة في كل نموذج،
أو في كل نموذج تصدق فيه البديهيات المختارة جميعًا.
فنستخرج أولًا الشروط التي تتيح ذلك، وهي شروط المنطق الجهوي
«النظامي». ويكفي فيه افتراض قاعدتين: القياس الاستثنائي والإلزام.
وأما البديهيات فجميع أمثلة إحلال القضايا الصادقة صدقًا تحليليًا،
ومعها ما يقتضيه المنطق المقصود من البديهيات الجهية.
وحتى إذا كان نظرنا في جميع النماذج، وجب إدخال كل أمثلة إحلال~\Ax{K}
\iftag{notprvDiamond}{}{و~$\Ax{Dual}$} في البديهيات؛ ومن هذا وحده
ينشأ أصغر منطق جهوي نظامي~$\Log K$.

\begin{defn}
قاعدة \emph{القياس الاستثنائي} هي مخطط الاستدلال
\begin{prooftree}
\AxiomC{$!A$}
\AxiomC{$!A \lif !B$}
\RightLabel{\MP}
\BinaryInfC{$!B$}
\end{prooftree}
نقول إن !!a{formula}~$!B$ \emph{تلزم عن}~!!{formula}s $!A$، $!C$
بالقياس الاستثنائي إذا وفقط إذا كان $!C \ident !A \lif !B$.
\end{defn}

\begin{defn}
قاعدة \emph{الإلزام} هي مخطط الاستدلال
\begin{prooftree}
\AxiomC{$!A$}
\RightLabel{\Nec}
\UnaryInfC{$\Box !A$}
\end{prooftree}
نقول إن !!{formula}~$!B$ تلزم عن !!{formula}s~$!A$ بالإلزام إذا وفقط
إذا كان $!B \ident \Box !A$.
\end{defn}

\begin{defn}
\emph{!!{derivation}} من مجموعة البديهيات~$\Sigma$ هو متتالية من
!!{formula}s $!B_{\OLMathNumeral{1}}$، $!B_{\OLMathNumeral{2}}$، \dots، $!B_{\OLId{latin-ordinary}{n}}$، بحيث يكون كل $!B_{\OLId{latin-ordinary}{i}}$ إما
\begin{enumerate}
\item مثال إحلال لقضية صادقة صدقا تحليليا، أو
\item مثال إحلال لـ!!a{formula} في~$\Sigma$، أو
\item لازما بالقياس الاستثنائي عن !!{formula}s اثنتين $!B_{\OLId{latin-ordinary}{j}}$، $!B_{\OLId{latin-ordinary}{k}}$
  حيث ${\OLId{latin-ordinary}{j}}$، ${\OLId{latin-ordinary}{k}} < {\OLId{latin-ordinary}{i}}$، أو
\item لازما بالإلزام عن !!a{formula}~$!B_{\OLId{latin-ordinary}{j}}$ حيث ${\OLId{latin-ordinary}{j}} < {\OLId{latin-ordinary}{i}}$.
\end{enumerate}
إذا وجد !!a{derivation} كهذا وكان $!B_{\OLId{latin-ordinary}{n}} \ident !A$، قلنا إن
$!A$ \emph{!!{derivable} من $\Sigma$}، ورمزنا إلى ذلك بـ$\Sigma \Proves !A$.
\end{defn}

وسيثبت أن الصيغ !!{derivable} بهذا التعريف تؤلف منطقًا جهويًا
نظاميًا، وأن كل عنصر !!{derivable} من نوع !!{formula} يصدق
في كل نموذج تصدق فيه البديهيات جميعًا. وهذه خاصية
\emph{السلامة} للـ!!{derivation}s. وأما العكس، وهو
\emph{التمام}، فإقامة برهانه أشق.

\end{document}
